% Corrected French for printed p.190; departures are append-only below.
\clearpage
\SourcePage{190}

\LayerTarget{tohoku:corollary:4.2.4:1}
\LayerTarget{tohoku:eq:4.2.10}
\[
\tag{4.2.10}
\mathbf{Ext}_O^n(A,B)=H^n(\mathbf{Hom}_O(\mathbf{L},B))
\]
d'où

\LayerTarget{tohoku:eq:4.2.10-bis}
\[
\tag{4.2.10 bis}
\mathbf{Ext}_O^n(A,B)(x)=
\mathrm{Ext}_{O(x)}^n(A(x),B(x)).
\]

\LayerTarget{tohoku:proof:corollary-4.2.4:1}
Il suffit en effet d'appliquer prop. 2.5.1.

Appliquons maintenant prop. 2.5.4, on trouve :

\LayerTarget{tohoku:corollary:4.2.4:2}
\noindent\textsc{Corollaire 2.} \textit{Sous les conditions précédentes, le
foncteur spectral du th. 4.2.1 est isomorphe au foncteur spectral
$\Pi\Gamma(\mathbf{Hom}_O(\mathbf{L},B))$, en particulier on a}
\[
\mathrm{Ext}_O^n(X;A,B)=
\mathbb{H}^n\Gamma(\mathbf{Hom}_O(\mathbf{L},B))
\]
\textit{(où le deuxième membre est le n.ième groupe de hyperhomologie de
$\Gamma$ par rapport au complexe $\mathbf{Hom}_O(\mathbf{L},B)$, cf. 2.4.).}

\LayerTarget{tohoku:construction:canonical-resolvent-spectral-sequence}
Comme nous disposons aussi d'un foncteur résolvant pour $\Gamma$, par exemple
le foncteur $\Gamma C(F)$ (où $C(F)$ est la \textit{résolution canonique} de
$F$, cf. 2.3), la suite spectrale
$\Pi\Gamma(\mathbf{Hom}_O(\mathbf{L},B))$ s'explicite, c'est la première suite
spectrale du bicomplexe $\Gamma C(\mathbf{Hom}_O(\mathbf{L},B))$ (où on a pris
comme premier degré celui qui provient de $C$). Or on a un isomorphisme
fonctoriel

\LayerTarget{tohoku:eq:4.2.11}
\[
\tag{4.2.11}
C(\mathbf{Hom}_O(\mathbf{L},B))=
\mathbf{Hom}_O(\mathbf{L},C(B)).
\]

En effet, on a pour tout $O$-Module $L$ un homomorphisme (fonctoriel)

\LayerTarget{tohoku:eq:4.2.12}
\[
\tag{4.2.12}
C(\mathbf{Hom}_O(L,B))\longrightarrow\mathbf{Hom}_O(L,C(B))
\]

\LayerTarget{tohoku:proof:equations-4.2.11-4.2.12}
(où au deuxième membre, $C(B)$ est considéré de façon évidente comme un
$O$-Module), défini par récurrence sur la dimension des composantes de $C$ à
partir de l'homomorphisme naturel
\[
C^0(\mathbf{Hom}_O(L,B))\longrightarrow\mathbf{Hom}_O(L,C^0(B)),
\]
déduit des homomorphismes naturels
$\mathbf{Hom}_O(L,B)(x)\longrightarrow
\mathrm{Hom}_{O(x)}(L(x),B(x))$. Ce dernier est un isomorphisme si $L$ est
pseudo-cohérent (prop. 4.1.1), d'où on conclut que dans ce cas, (4.2.12) est un
isomorphisme, ce qui établit en particulier la formule 4.2.11. On en conclut
$\Gamma C(\mathbf{Hom}_O(\mathbf{L},B))=
\mathbf{Hom}_O(\mathbf{L},C(B))$, d'où le

\LayerTarget{tohoku:corollary:4.2.4:3}
\noindent\textsc{Corollaire 3.} \textit{$\mathbf{L}$ étant comme dans la
proposition 4.2.4, le foncteur spectral du théorème 4.2.1 est donné par la
première suite spectrale
$\Gamma\mathbf{Hom}_O(\mathbf{L},C(B))$ du bicomplexe
$\mathbf{Hom}_O(\mathbf{L},C(B))$, où $C(B)$ est la résolution canonique
(cf. 2.3) du faisceau abélien $B$ (dont la graduation est prise comme premier
degré du bicomplexe). En particulier, on a}
\[
\mathrm{Ext}_O^n(X;A,B)=
H^n(\mathbf{Hom}_O(\mathbf{L},C(B))).
\]

\LayerTarget{tohoku:remark:ext-with-supports}
\noindent\textsc{Remarques.} 1) On peut aussi définir les foncteurs
$\mathrm{Ext}_{O,\varphi}^p(X;A,B)$ comme les foncteurs dérivés droits, par
rapport à $B$, du foncteur
$\mathrm{Hom}_{O,\varphi}(A,B)=
\Gamma_\varphi(\mathbf{Hom}_O(A,B))$ ($\varphi$ étant un antifiltre de parties
fermées de $X$). Les développements précédents sont encore valables. Nous
n'aurons pas à nous en servir.

\LayerTarget{tohoku:remark:finiteness-in-theorem-4.2.2}
2) Il est essentiel dans le théorème 4.2.2 que $A$ satisfasse à une condition
de finitude. Si par exemple $A=O^{(I)}$ où $I$ est infini, alors le deuxième
membre de (4.2.8) est nul pour tout $x$ et tout $p>0$, cependant on n'aura
\textit{pas} en général $\mathbf{Ext}_O^1(A,B)=0$, car le foncteur
$B\longmapsto\mathbf{Hom}_O(O^{(I)},B)=B^I$ n'est en général pas un foncteur
exact.

\TohokuContinuation{191}
